% Options for packages loaded elsewhere
\PassOptionsToPackage{unicode}{hyperref}
\PassOptionsToPackage{hyphens}{url}
\PassOptionsToPackage{dvipsnames,svgnames,x11names}{xcolor}
%
\documentclass[
  10pt,
]{article}
\usepackage{amsmath,amssymb}
\usepackage{iftex}
\ifPDFTeX
  \usepackage[T1]{fontenc}
  \usepackage[utf8]{inputenc}
  \usepackage{textcomp} % provide euro and other symbols
\else % if luatex or xetex
  \usepackage{unicode-math} % this also loads fontspec
  \defaultfontfeatures{Scale=MatchLowercase}
  \defaultfontfeatures[\rmfamily]{Ligatures=TeX,Scale=1}
\fi
\usepackage{lmodern}
\ifPDFTeX\else
  % xetex/luatex font selection
\fi
% Use upquote if available, for straight quotes in verbatim environments
\IfFileExists{upquote.sty}{\usepackage{upquote}}{}
\IfFileExists{microtype.sty}{% use microtype if available
  \usepackage[]{microtype}
  \UseMicrotypeSet[protrusion]{basicmath} % disable protrusion for tt fonts
}{}
\makeatletter
\@ifundefined{KOMAClassName}{% if non-KOMA class
  \IfFileExists{parskip.sty}{%
    \usepackage{parskip}
  }{% else
    \setlength{\parindent}{0pt}
    \setlength{\parskip}{6pt plus 2pt minus 1pt}}
}{% if KOMA class
  \KOMAoptions{parskip=half}}
\makeatother
\usepackage{xcolor}
\usepackage[margin=1in]{geometry}
\usepackage{color}
\usepackage{fancyvrb}
\newcommand{\VerbBar}{|}
\newcommand{\VERB}{\Verb[commandchars=\\\{\}]}
\DefineVerbatimEnvironment{Highlighting}{Verbatim}{commandchars=\\\{\}}
% Add ',fontsize=\small' for more characters per line
\newenvironment{Shaded}{}{}
\newcommand{\AlertTok}[1]{\textcolor[rgb]{1.00,0.00,0.00}{\textbf{#1}}}
\newcommand{\AnnotationTok}[1]{\textcolor[rgb]{0.38,0.63,0.69}{\textbf{\textit{#1}}}}
\newcommand{\AttributeTok}[1]{\textcolor[rgb]{0.49,0.56,0.16}{#1}}
\newcommand{\BaseNTok}[1]{\textcolor[rgb]{0.25,0.63,0.44}{#1}}
\newcommand{\BuiltInTok}[1]{\textcolor[rgb]{0.00,0.50,0.00}{#1}}
\newcommand{\CharTok}[1]{\textcolor[rgb]{0.25,0.44,0.63}{#1}}
\newcommand{\CommentTok}[1]{\textcolor[rgb]{0.38,0.63,0.69}{\textit{#1}}}
\newcommand{\CommentVarTok}[1]{\textcolor[rgb]{0.38,0.63,0.69}{\textbf{\textit{#1}}}}
\newcommand{\ConstantTok}[1]{\textcolor[rgb]{0.53,0.00,0.00}{#1}}
\newcommand{\ControlFlowTok}[1]{\textcolor[rgb]{0.00,0.44,0.13}{\textbf{#1}}}
\newcommand{\DataTypeTok}[1]{\textcolor[rgb]{0.56,0.13,0.00}{#1}}
\newcommand{\DecValTok}[1]{\textcolor[rgb]{0.25,0.63,0.44}{#1}}
\newcommand{\DocumentationTok}[1]{\textcolor[rgb]{0.73,0.13,0.13}{\textit{#1}}}
\newcommand{\ErrorTok}[1]{\textcolor[rgb]{1.00,0.00,0.00}{\textbf{#1}}}
\newcommand{\ExtensionTok}[1]{#1}
\newcommand{\FloatTok}[1]{\textcolor[rgb]{0.25,0.63,0.44}{#1}}
\newcommand{\FunctionTok}[1]{\textcolor[rgb]{0.02,0.16,0.49}{#1}}
\newcommand{\ImportTok}[1]{\textcolor[rgb]{0.00,0.50,0.00}{\textbf{#1}}}
\newcommand{\InformationTok}[1]{\textcolor[rgb]{0.38,0.63,0.69}{\textbf{\textit{#1}}}}
\newcommand{\KeywordTok}[1]{\textcolor[rgb]{0.00,0.44,0.13}{\textbf{#1}}}
\newcommand{\NormalTok}[1]{#1}
\newcommand{\OperatorTok}[1]{\textcolor[rgb]{0.40,0.40,0.40}{#1}}
\newcommand{\OtherTok}[1]{\textcolor[rgb]{0.00,0.44,0.13}{#1}}
\newcommand{\PreprocessorTok}[1]{\textcolor[rgb]{0.74,0.48,0.00}{#1}}
\newcommand{\RegionMarkerTok}[1]{#1}
\newcommand{\SpecialCharTok}[1]{\textcolor[rgb]{0.25,0.44,0.63}{#1}}
\newcommand{\SpecialStringTok}[1]{\textcolor[rgb]{0.73,0.40,0.53}{#1}}
\newcommand{\StringTok}[1]{\textcolor[rgb]{0.25,0.44,0.63}{#1}}
\newcommand{\VariableTok}[1]{\textcolor[rgb]{0.10,0.09,0.49}{#1}}
\newcommand{\VerbatimStringTok}[1]{\textcolor[rgb]{0.25,0.44,0.63}{#1}}
\newcommand{\WarningTok}[1]{\textcolor[rgb]{0.38,0.63,0.69}{\textbf{\textit{#1}}}}
\setlength{\emergencystretch}{3em} % prevent overfull lines
\providecommand{\tightlist}{%
  \setlength{\itemsep}{0pt}\setlength{\parskip}{0pt}}
\setcounter{secnumdepth}{-\maxdimen} % remove section numbering
\ifLuaTeX
  \usepackage{selnolig}  % disable illegal ligatures
\fi
\usepackage{bookmark}
\IfFileExists{xurl.sty}{\usepackage{xurl}}{} % add URL line breaks if available
\urlstyle{same}
\hypersetup{
  colorlinks=true,
  linkcolor={Maroon},
  filecolor={Maroon},
  citecolor={Blue},
  urlcolor={blue},
  pdfcreator={LaTeX via pandoc}}

\author{}
\date{}

\begin{document}

\section{BA1-T2 - Split-or-collapse pressure test of the two-family
candidate}\label{ba1-t2---split-or-collapse-pressure-test-of-the-two-family-candidate}

\textbf{Revision:} R1

\textbf{Status:} COMPLETED / PROVISIONAL PASS WITH MINIMALITY REFINEMENT
/ BA1 NOT CLOSED

\textbf{Repository baseline reviewed:}
\texttt{f05fbb7b253392e158a1062df2614b177c13d43e}

\textbf{Phase:} BA1 - Minimal BAE ontology

\textbf{BA0:} CLOSED

\textbf{BA2:} NOT STARTED

\subsection{1. Question and test
discipline}\label{question-and-test-discipline}

BA1-T2 asks only:

\begin{quote}
Does the BA1-T1 candidate \texttt{BAReferent\ +\ BAProposition} survive
deliberate pressure to split recurring semantic kinds into dedicated
first-class types, and does it also survive the opposite pressure to
collapse both families into one undifferentiated analytical element?
\end{quote}

The test does not design the BA2 relation/action vocabulary, BA3
provenance schema, BA4 views, formal STRIDE overlay, Finding model or
ThreatForge implementation.

The governing criterion remains:

\begin{Shaded}
\begin{Highlighting}[]
\NormalTok{independent semantic identity}
\NormalTok{        +}
\NormalTok{reusable subtype{-}specific invariants}
\NormalTok{        {-}\textgreater{} possible first{-}class split}

\NormalTok{recurring usefulness alone}
\NormalTok{        !=}
\NormalTok{first{-}class metaclass}
\end{Highlighting}
\end{Shaded}

A pressure case may justify an independently identifiable
\texttt{BAReferent} without justifying a dedicated subtype such as
\texttt{Behavior}, \texttt{Information}, \texttt{State},
\texttt{Contract} or \texttt{Store}.

\subsection{2. Starting candidate from
BA1-T1}\label{starting-candidate-from-ba1-t1}

BA1-T1 proposed two identity families:

\begin{Shaded}
\begin{Highlighting}[]
\NormalTok{BAReferent}
\NormalTok{BAProposition}
\end{Highlighting}
\end{Shaded}

\texttt{BAReferent} covered independently identifiable units of shared
project meaning. \texttt{BAProposition} covered independently
identifiable methodology-neutral analytical assertions with
origin/change/diagnostic pressure.

BA1-T1 deliberately left domain splits open and allowed a proposition to
be targeted by another proposition. BA1-T2 treats both choices as
falsifiable rather than inherited truth.

\subsection{3. Pressure A - behavior/event
identity}\label{pressure-a---behaviorevent-identity}

\subsubsection{3.1 Facial-access delivery}\label{facial-access-delivery}

\texttt{FR-3.4} requires delivery of \texttt{RecognitionCapture} from
\texttt{CameraSubsystem} to \texttt{RecognitionProcessor}, correlated to
the correct \texttt{RecognitionRequest}, with incomplete delivery not
represented as success. Specialized requirements independently constrain
confidentiality, integrity and authorized provenance during that same
delivery meaning.

The pressure question is whether this requires a dedicated
\texttt{Behavior} or \texttt{DataFlow} metaclass.

It does not.

The delivery meaning can have independent analytical identity as a
\texttt{BAReferent} when multiple propositions need to talk about the
same behavior. Separate propositions can state its source, destination,
information, correlation and failure semantics, together with the
specialized security constraints that apply to it.

This preserves behavior identity without asserting that every
behavior/event belongs to a dedicated root metaclass.

\subsubsection{3.2 Physical handoff
milestone}\label{physical-handoff-milestone}

The order-fulfillment corpus gives stronger pressure. The physical
handoff milestone is defined in the fulfillment branch and reused by
inventory stock issue and payment capture. Reuse across branches makes
the milestone a natural independently identifiable semantic object.

Again, the evidence forces \textbf{referent identity}, not a universal
\texttt{Event} metaclass. The milestone can be a \texttt{BAReferent};
propositions state who confirms it, when it occurs, and which
stock/payment behavior depends on it.

\subsubsection{3.3 Disposition}\label{disposition}

\begin{Shaded}
\begin{Highlighting}[]
\NormalTok{Behavior / Event dedicated first{-}class split: NOT FORCED}
\NormalTok{Behavior/Event referent identity when independently reused: SUPPORTED}
\end{Highlighting}
\end{Shaded}

The old historical \texttt{data\_flow} focus is not revived.

\subsection{4. Pressure B - information/resource identity and
lifecycle}\label{pressure-b---informationresource-identity-and-lifecycle}

The two corpora contain strongly reusable informational/resource-like
meanings:

\begin{itemize}
\tightlist
\item
  \texttt{RecognitionCapture};
\item
  \texttt{RecognitionRequest};
\item
  \texttt{OrderEvaluation};
\item
  \texttt{Order};
\item
  \texttt{Reservation};
\item
  \texttt{ReservationResult};
\item
  \texttt{PaymentAuthorizationResult};
\item
  \texttt{FulfillmentTask}.
\end{itemize}

Some carry lifecycle or correlation semantics. \texttt{Reservation}, for
example, is bound to one \texttt{OrderEvaluation}, has an ACTIVE
lifecycle state, can be released/expired/consumed and participates in
idempotent handling.

Those facts require stable referent identity and propositions about
lifecycle/correlation. They do not establish one reusable set of
invariants that all information/resource referents share and that
generic referents cannot preserve.

A subtype may later become useful, but current evidence does not make it
necessary for correctness or method-neutral reuse.

\subsubsection{Disposition}\label{disposition-1}

\begin{Shaded}
\begin{Highlighting}[]
\NormalTok{Information / Resource dedicated first{-}class split: NOT FORCED}
\NormalTok{Lifecycle{-}bearing informational referents: representable as BAReferent + propositions}
\end{Highlighting}
\end{Shaded}

\subsection{5. Pressure C - responsibility, externality and
boundaries}\label{pressure-c---responsibility-externality-and-boundaries}

Facial access distinguishes project endpoints from underlying external
transport responsibility. Order fulfillment distinguishes internal
project-owned inventory authority from an alternative external WMS and
places payment/carrier providers behind governed contracts.

The pressure question is whether this forces dedicated
\texttt{Participant}, \texttt{Component}, \texttt{Capability},
\texttt{Boundary} or \texttt{ExternalSystem} roots.

It does not.

The independently identifiable parties/capabilities are referents.
Project ownership, externality, service consumption and responsibility
placement are propositions or roles over those referents. A named
responsibility boundary may itself receive referent identity if the
project needs to reuse that boundary explicitly, but no subtype-specific
invariant is yet forced.

The critical semantic point is preserved without type proliferation:

\begin{Shaded}
\begin{Highlighting}[]
\NormalTok{service/capability consumption}
\NormalTok{        !=}
\NormalTok{ownership transfer}
\end{Highlighting}
\end{Shaded}

\subsubsection{Disposition}\label{disposition-2}

\begin{Shaded}
\begin{Highlighting}[]
\NormalTok{Participant / Component / Capability split: NOT FORCED}
\NormalTok{Boundary / ExternalSystem split: NOT FORCED}
\NormalTok{Responsibility and externality: proposition/role pressure supported}
\end{Highlighting}
\end{Shaded}

\subsection{6. Pressure D - Store and Contract
candidates}\label{pressure-d---store-and-contract-candidates}

The order corpus names the following candidates:

\begin{itemize}
\tightlist
\item
  stores: \texttt{InventoryLedger}, \texttt{ReservationStore};
\item
  contracts: \texttt{PaymentProviderContract}, \texttt{CarrierContract}.
\end{itemize}

All require independent identity if they participate in analysis. But
the corpus does not establish universal subtype invariants that cannot
be stated over generic referents:

\begin{itemize}
\tightlist
\item
  a store may be project-owned or merely a candidate realization;
\item
  a contract may define a governed interaction boundary while
  provider-specific raw states remain outside project semantics;
\item
  replacement of a provider/contract does not imply that all contracts
  or stores require a dedicated root type for every consumer.
\end{itemize}

Their semantic kind must not be lost, but classification is not
equivalent to a first-class metaclass.

\subsubsection{Disposition}\label{disposition-3}

\begin{Shaded}
\begin{Highlighting}[]
\NormalTok{Store dedicated first{-}class split: NOT FORCED}
\NormalTok{Contract dedicated first{-}class split: NOT FORCED}
\end{Highlighting}
\end{Shaded}

\subsection{7. Pressure E - state, context and indeterminate
outcomes}\label{pressure-e---state-context-and-indeterminate-outcomes}

The order corpus uses values/states such as:

\begin{itemize}
\tightlist
\item
  \texttt{Order=ACCEPTED};
\item
  payment \texttt{authorized}, \texttt{declined},
  \texttt{indeterminate};
\item
  reservation ACTIVE and terminal lifecycle states;
\item
  availability/result states such as
  sufficient/insufficient/indeterminate.
\end{itemize}

Current evidence supports three honest representations depending on
semantic need:

\begin{enumerate}
\def\labelenumi{\arabic{enumi}.}
\tightlist
\item
  scalar/enumerated value or qualifier when only local value semantics
  are needed;
\item
  proposition when a state condition or transition must be asserted;
\item
  \texttt{BAReferent} when a named state/mode/context has independent
  reusable identity across propositions.
\end{enumerate}

No current case forces a universal \texttt{State} or \texttt{Context}
metaclass.

\subsubsection{Disposition}\label{disposition-4}

\begin{Shaded}
\begin{Highlighting}[]
\NormalTok{State / Mode / Context dedicated first{-}class split: NOT FORCED}
\end{Highlighting}
\end{Shaded}

\subsection{8. Pressure F -
proposition-as-target}\label{pressure-f---proposition-as-target}

BA1-T1 allowed a \texttt{BAProposition} to be targeted by another
proposition. The pressure cases show that this is not required by the
minimal ontology and can blur two different meanings:

\begin{Shaded}
\begin{Highlighting}[]
\NormalTok{project meaning being constrained}
\NormalTok{        !=}
\NormalTok{analytical assertion record describing that meaning}
\end{Highlighting}
\end{Shaded}

For example, \texttt{SR-3.4-C}, \texttt{SR-3.4-I} and \texttt{SR-3.4-P}
constrain the delivery behavior itself. They do not semantically
constrain the existence of the analytical proposition record used to
describe the delivery.

The cleaner minimal rule is:

\begin{itemize}
\tightlist
\item
  if project meaning needs independent reuse/qualification, give that
  meaning \texttt{BAReferent} identity;
\item
  \texttt{BAProposition} asserts methodology-neutral meaning about
  referents;
\item
  provenance, diagnostic and change-state metadata may target
  proposition identity in BA3, but that is analytical metadata rather
  than a project-semantic proposition.
\end{itemize}

BA1-T2 therefore removes proposition-as-project-target as a required
feature of the minimal candidate.

\subsubsection{Disposition}\label{disposition-5}

\begin{Shaded}
\begin{Highlighting}[]
\NormalTok{BAProposition as target of project{-}semantic BAProposition: NOT REQUIRED / REMOVED FROM MINIMAL CANDIDATE}
\NormalTok{Proposition identity for provenance/change/diagnostic handling: STILL REQUIRED}
\end{Highlighting}
\end{Shaded}

This is a reduction of R1, not a new type.

\subsection{9. Pressure G - bounded method-consumer
projection}\label{pressure-g---bounded-method-consumer-projection}

The BA0-T2 bounded STRIDE transfer consumer requires a projection
containing:

\begin{itemize}
\tightlist
\item
  behavior;
\item
  source;
\item
  destination;
\item
  information;
\item
  correlation;
\item
  responsibility/externality;
\item
  current realization;
\item
  failure semantics;
\item
  confidentiality/integrity/authorized-origin constraints.
\end{itemize}

The two-family ontology can supply the needed identities without
importing STRIDE or DFD types:

\begin{itemize}
\tightlist
\item
  independently reused project meanings are \texttt{BAReferent}
  identities;
\item
  methodology-neutral assertions connecting and qualifying those
  referents are \texttt{BAProposition} identities;
\item
  BA2 will define the minimal shared relation/action roles needed to
  make those propositions mechanically selectable;
\item
  the STRIDE consumer adds its own categories and dispositions outside
  Base Analysis.
\end{itemize}

The consumer therefore does not need a dedicated \texttt{Actor},
\texttt{Process}, \texttt{DataStore}, \texttt{DataFlow} or
\texttt{TrustBoundary} root in BA1.

This is an ontology-level pass, not yet a proof of the final BA2/BA4
material projection contract.

\subsubsection{Disposition}\label{disposition-6}

\begin{Shaded}
\begin{Highlighting}[]
\NormalTok{Bounded method consumer: PASS AT BA1 IDENTITY{-}ONTOLOGY LEVEL}
\NormalTok{Method{-}specific type contamination: NOT REQUIRED}
\end{Highlighting}
\end{Shaded}

\subsection{10. Collapse pressure - can the two families become
one?}\label{collapse-pressure---can-the-two-families-become-one}

\subsubsection{10.1 Strong collapse
candidate}\label{strong-collapse-candidate}

The genuine collapse hypothesis is:

\begin{Shaded}
\begin{Highlighting}[]
\NormalTok{BAElement}
\end{Highlighting}
\end{Shaded}

with no normative semantic distinction between project referent identity
and analytical assertion identity.

This fails.

The same referent can remain stable while propositions about it are
introduced, revised, superseded, invalidated or diagnosed. In facial
access, \texttt{CameraSubsystem} can persist while delivery/realization
statements change under mutations. In order fulfillment, one
\texttt{OrderEvaluation} participates in independently evolving
propositions about reservation, payment, compensation and commitment.

Claim-level origin/change/diagnostic handling therefore cannot be
reduced to referent identity without losing a real semantic distinction.

\subsubsection{10.2 Physical single-class implementation is not semantic
collapse}\label{physical-single-class-implementation-is-not-semantic-collapse}

A storage implementation could materialize both families in one
table/class with a mandatory discriminator. That does \textbf{not}
falsify the two-family ontology: the semantic distinction would still
exist as a normative partition with different invariants.

BA1 concerns semantic identity families, not the number of programming
classes or database tables.

\subsubsection{Disposition}\label{disposition-7}

\begin{Shaded}
\begin{Highlighting}[]
\NormalTok{Semantic collapse to one undifferentiated family: REJECTED}
\NormalTok{Single physical implementation class/table: REPRESENTATION CHOICE / IRRELEVANT TO BA1}
\end{Highlighting}
\end{Shaded}

\subsection{11. Split-pressure summary}\label{split-pressure-summary}

\begin{itemize}
\tightlist
\item
  Behavior/Event: \textbf{NOT FORCED} as dedicated type; independent
  event/behavior identity can be \texttt{BAReferent}.
\item
  Information/Resource: \textbf{NOT FORCED}.
\item
  Participant/Component/Capability: \textbf{NOT FORCED}.
\item
  Boundary/Externality: \textbf{NOT FORCED}.
\item
  Store: \textbf{NOT FORCED}.
\item
  Contract: \textbf{NOT FORCED}.
\item
  State/Mode/Context: \textbf{NOT FORCED}.
\item
  Method-specific DFD/STRIDE roots: \textbf{REJECTED AS BA1 NECESSITY}.
\item
  Proposition-as-project-target: \textbf{NOT REQUIRED / REMOVED FROM
  MINIMAL CANDIDATE}.
\item
  Semantic one-family collapse: \textbf{REJECTED}.
\end{itemize}

No tested case reaches \texttt{FORCED\ CANDIDATE\ SPLIT}.

\subsection{12. Refined two-family
candidate}\label{refined-two-family-candidate}

BA1-T2 therefore retains exactly two candidate identity families, with a
sharper boundary.

\subsubsection{\texorpdfstring{BA1-C1 -
\texttt{BAReferent}}{BA1-C1 - BAReferent}}\label{ba1-c1---bareferent}

\textbf{Status:} STRONG CANDIDATE / NOT ACCEPTED

An independently identifiable unit of methodology-neutral shared project
meaning that must be reusable across one or more propositions,
projections or governed baselines.

Its semantic kind may matter and must not be lost, but BA1 does not yet
require dedicated metaclasses for participant, information,
behavior/event, store, contract, boundary, state or similar kinds.

When behavior/event/state/context meaning needs independent
project-semantic reuse, it can itself be represented as a
\texttt{BAReferent}.

\subsubsection{\texorpdfstring{BA1-C2 -
\texttt{BAProposition}}{BA1-C2 - BAProposition}}\label{ba1-c2---baproposition}

\textbf{Status:} STRONG CANDIDATE / NOT ACCEPTED

An independently identifiable methodology-neutral analytical assertion
about one or more \texttt{BAReferent} identities, with enough identity
for origin/provenance, diagnosis, reuse and change disposition.

A proposition is not itself the project meaning it describes.
Project-semantic qualification should target referents, while
proposition provenance/diagnostic state is handled as analytical
metadata in BA3.

The exact proposition structure, predicates, roles, n-ary representation
and action vocabulary remain BA2 questions.

\subsection{13. What BA1-T2 does not
decide}\label{what-ba1-t2-does-not-decide}

BA1-T2 does not decide:

\begin{itemize}
\tightlist
\item
  whether \texttt{BAReferent} and \texttt{BAProposition} share an
  abstract implementation superclass;
\item
  the canonical semantic-kind/classification vocabulary;
\item
  relation cardinalities or n-ary proposition encoding;
\item
  action/predicate vocabulary;
\item
  identity/equivalence/lifecycle mechanics;
\item
  source locators and provenance records;
\item
  diagnostic record schema;
\item
  view/projection contracts;
\item
  lexical controlled vocabulary;
\item
  AnalysisRecord, Finding or method-specific schemas.
\end{itemize}

These remain in BA2-BA6 or the later analysis envelope.

\subsection{14. Falsification
assessment}\label{falsification-assessment}

The two-family candidate would have failed if any tested recurring kind
required independent subtype-specific invariants that generic referent
identity could not preserve, or if the bounded method consumer needed
method-specific roots in the shared core.

No such counterexample was found in the two current corpora.

The candidate would also have been reducible if proposition identity
could be eliminated without losing independent origin/change/diagnostic
handling. The collapse test failed.

One R1 feature was weakened: proposition-as-project-target is
unnecessary and risks confusing assertion records with project meaning.
Removing it improves minimality.

\subsection{15. Trial disposition}\label{trial-disposition}

\begin{Shaded}
\begin{Highlighting}[]
\NormalTok{BA1{-}T2                         COMPLETED / PROVISIONAL PASS WITH REFINEMENT}
\NormalTok{BAReferent                     STRONG CANDIDATE / NOT ACCEPTED}
\NormalTok{BAProposition                  STRONG CANDIDATE / NOT ACCEPTED}
\NormalTok{Dedicated domain{-}type splits   NOT FORCED BY TESTED EVIDENCE}
\NormalTok{Semantic one{-}family collapse   REJECTED}
\NormalTok{BA1                            OPEN / READY FOR DISTINCT CLOSURE REVIEW}
\NormalTok{BA2                            NOT STARTED}
\end{Highlighting}
\end{Shaded}

BA1 is \textbf{not closed by this trial}. The evidence is now stable
enough to perform a distinct closure review rather than inventing
another ontology experiment for reassurance.

\subsection{16. Next authorized
microstep}\label{next-authorized-microstep}

Execute only:

\begin{quote}
\textbf{BA1-T3 - minimal BAE ontology closure review.}
\end{quote}

BA1-T3 must test whether the refined two-family boundary is necessary,
sufficient and correctly phased, and whether every remaining unresolved
design question genuinely belongs to BA2-BA6.

If no material counterexample remains, BA1-T3 may close BA1 and
authorize BA2. It must not begin relation/action vocabulary design
inside the closure review.

\end{document}
